\documentclass[french]{book}
\usepackage{teach}
%\usepackage{verbatim}
\usepackage{graphicx}
\graphicspath{{Figures/}}
\usepackage{amsmath,amsfonts,amssymb}
\usepackage{tikz,tkz-tab}
\usepackage{multicol}
\redefinecolor{thm}{title}{bgcolor}{alizarin}
\redefinecolor{prop}{title}{bgcolor}{meatbrown}
\redefinecolor{defn}{title}{bgcolor}{olivine}
\definecolor{alizarin}{rgb}{0.82, 0.1, 0.26}
\definecolor{meatbrown}{rgb}{0.9, 0.72, 0.23}
\definecolor{olivine}{rgb}{0.6, 0.73, 0.45}
%\usepackage{mathrsfs}
%%
\newcommand{\R}{\mathbb {R}}
%\newcommand{\C}{\mathds {C}}
\newcommand{\N}{\mathbb {N}}
\newcommand{\Z}{\mathbb {Z}}
\newcommand{\Q}{\mathbb {Q}}
\newcommand{\D}{\mathbb {D}}
%%
\newcommand{\se}{\geq}
\newcommand{\ie}{\leq}
\newcommand{\tm}{\times}
%\newcommand{\ell}{l}
%\newcommand{\ell'}{l'}
%%
\newcommand{\limn}{\lim_{n\rightarrow +\infty}}
\newcommand{\limo}[1][x]{\ds\lim_{#1\rightarrow 0}}
\newcommand{\limite}[2][x]{\ds\lim_{#1\rightarrow #2}}
\newcommand{\limpinf}[1][x]{\ds\lim_{#1\rightarrow +\infty}}
\newcommand{\limminf}[1][x]{\ds\lim_{#1\rightarrow -\infty}}
\newcommand{\limiteg}[2][x]{\ds\lim_{#1\xrightarrow{<}#2}}
\newcommand{\limited}[2][x]{\ds\lim_{#1\xrightarrow{>}#2}}
%%
\newcommand{\vect}[1]{\overrightarrow{#1}}
\newcommand{\oij}{(\textit{O},{\overrightarrow{i}},{\overrightarrow{j}})}
%
\newcommand{\cp}[2]{%
        \begin{pmatrix}
        #1\\
        #2
        \end{pmatrix}%
        }
%
\newcommand{\calb}{\mathscr{B}}
\newcommand{\calc}{\mathscr{C}}
\newcommand{\cald}{\mathscr{D}}
\newcommand{\cale}{\mathscr{E}}
\newcommand{\calf}{\mathscr{F}}
\newcommand{\calp}{\mathscr{P}}
\newcommand{\cals}{\mathscr{S}}
\newcommand{\calt}{\mathscr{T}}
%
\newcommand{\suite}[1][u]{\left( #1_{n} \right)_{n\in\N}}
\newcommand{\suitep}[1][u]{\left( #1_{n} \right)_{n\in\N^{*}}}


\newcommand{\poly}[1][a]{#1_0+#1_1x+\cdots+#1_nx^n}
\newcommand{\defi}[1]{\begin{defn} #1 \end{defn}}

\usepackage{dashundergaps}
\newcommand{\mgap}[1]{\text{\gap*[blank]{\ensuremath{#1}}}}

\begin{document}

\chapter{Polynômes du second et du troisième degré}

\section{Rappel}
\begin{prop}
Pour tous les nombres réels $a$, $b$ et $k$ :
\begin{multicols}{2}
$$k(a+b)=ka+kb$$
$$(a+b)(c+d)=ac+ad+bc+bd$$
$$(a+b)^2=a^2+2ab+b^2$$
$$(a-b)^2=a^2-2ab+b^2$$
$$(a-b)(a+b)=a^2-b^2$$
\end{multicols}
\end{prop}
\section{Définitions et propriétés}
\begin{defn}
On appelle \gap*{fonction polynôme du second degré} toute fonction $f$ définie sur $\mathbb{R}$ et qui s'écrit $f(x)=ax^2+bx+c$ où $a$, $b$ et $c$ sont des réels fixés et $a\neq 0$.
\end{defn}

\begin{exemple}
la fonction $f:x\mapsto -x^{3}-5x^{2}+7x-1$  n'est pas une fonction polynôme du second degré. la fonction $g:x\mapsto 3x-x^{2}$ est une fonction polynôme du second degré avec $a=-1$ ,$b=3$ et $c=0$.
\end{exemple}

\subsection{Racines d'un polynôme}
\begin{defn}
Soit $f$ une fonction polynôme du second degré. On appelle \gap*{racine} de $f$ toute solution de l'équation $f(x)=0$.
\vspace{5mm}
\end{defn}

\begin{exemple}
0 est une racine de $x\mapsto 3x-x^{2}$ car $3\times0-0^2=0$. Attention, certains polynômes n'admettent pas de racine.
\end{exemple}

\section{Equation du type $x^2=c$, avec $c$ strictement positif}

\begin{thm}
Soit $c$ un nombre \gap*{strictement positif} alors l'équation du type $x^2=c$ admet deux solutions notées $x_1$, $x_2$  et telles que $x_1=\sqrt{c}$ et $x_2=-\sqrt{c}$.
\end{thm}

\begin{exemple}
$x^2=2$ admet deux solutions $x_1=\sqrt{2}$ et $x_2=-\sqrt{2}$.
\end{exemple}

\section{Forme factorisée et forme développée d'un polynôme du second degré}

Un polynôme du second degré peut être écrit sous deux formes différentes : \gap*{la forme développée} et \gap*{potentiellement sous forme factorisée}.

\begin{defn}
La \gap*{forme développée} d'un polynôme du second degré est donnée par $ax^2 + bx + c$, où $a, b$ et $c$ sont des \gap*{coefficients}.
\end{defn}

\defi{Un polynôme du second degré peut être écrit sous \gap*{forme factorisée} sous la forme $a(x - x_1)(x - x_2)$, où $x_1$ et $x_2$ sont les \gap*{racines du polynôme}.}

\begin{exemple}
Soit le polynôme du second degré $x^2 - 5x + 6$. Nous pouvons le factoriser en $(x - 2)(x - 3)$. Il est donc possible de passer de la forme développée à la forme factorisée et vice versa.

Attention, un polynôme qui \gap*{ne possède pas de racine} ne possède pas non plus de \gap*{forme factorisée}.
\end{exemple}


\subsection{Relations de Viète pour les polynômes du second degré}

\begin{prop}
Soit $ax^2+bx+c$ un polynôme du second degré tel qu'il existe admet deux racines $x_{1}$ et $x_{2}$ alors 
$\left\lbrace 
\begin{array}[c]{l}
x_{1}+x_{2}=-\dfrac{b}{a} \\
x_{1}x_{2}=\dfrac{c}{a}
\end{array}
\right.$
\end{prop}

\newpage
\begin{exemple}
Soit le polynôme suivant : $f(x) = 2x^2 - 15x + 25$. On sait que l'une des racines de ce polynôme est 5. Utilisez les relations de Viète pour déterminer l'autre racine. \\

Les relations de Viète sont des formules qui établissent des relations entre les coefficients d'un polynôme et ses racines. Pour un polynôme de degré 2 de la forme $ax^2 + bx + c$, ces relations sont les suivantes :\\

\begin{itemize}
\item La somme des racines est égale à $-\dfrac{b}{a}$
\item Le produit des racines est égal à $\dfrac{c}{a}$
\end{itemize}

Dans notre cas, le polynôme $f(x) = 2x^2 - 15x + 25$ est de degré 2 et de la forme $ax^2 + bx + c$ avec $a=2$, $b=-15$ et $c=25$. Nous pouvons donc utiliser les relations de Viète pour trouver les racines de ce polynôme.\\

Calcul de la somme des racines
La somme des racines est égale à $-\dfrac{b}{a}$. Dans notre cas, cela donne :

$$-\dfrac{b}{a} = \dfrac{15}{2}$$

Calcul du produit des racines
Le produit des racines est égal à $\dfrac{c}{a}$. Dans notre cas, cela donne :

$$\dfrac{c}{a} = \dfrac{25}{2}$$

Nous savons maintenant que la somme des racines est $\dfrac{15}{2}$ et que le produit des racines est $\dfrac{25}{2}$. Nous pouvons utiliser ces informations pour trouver l'autre racine du polynôme. Pour cela, nous devons trouver un nombre qui, ajouté à 5, donne une somme égale à $\dfrac{15}{2}$ et un produit égal à $\dfrac{25}{2}$.

Nous pouvons résoudre ce système d'équations:

$$x + 5 = \frac{15}{2}$$
$$5x = \frac{25}{2}$$

En utilisant une des deux équations, au choix, pour exprimer $x$, on obtient $x = \dfrac{5}{2}$
La deuxième racine du polynôme est donc $ \dfrac{5}{2}$.

Vérification :
Nous pouvons vérifier que ces valeurs sont bien les racines du polynôme en utilisant la formule de la factorisation du polynôme. En effet, certains polynômes de degré 2 peuvent s'écrire sous la forme $a(x-r_1)(x-r_2)$, où $r_1$ et $r_2$ sont les racines du polynôme.

Dans notre cas, le polynôme $f(x) = 2x^2 - 15x + 25$ peut s'écrire sous la forme :
$$f(x) = 2(x - 5)(x - \dfrac{5}{2})$$
Nous pouvons donc vérifier que les racines du polynôme sont bien 5 et $\dfrac{5}{2}$.
\end{exemple}


\section{Représentation graphique des polynômes du type $x\mapsto ax^2$ et $x\mapsto ax^2+b$}

\begin{exemple}

Représentation graphique de $w$ définie par $w(x)=0,07x^2$ sur $[-1,1;1,1]$.

\begin{center}
\includegraphics[scale=0.9]{courbeax2b.3}
\end{center}

\end{exemple}

\begin{exemple}


Représentation graphique de $z$ définie par $z(x)=-3x^2+\dfrac{3}{2}$ sur $[-1,1;1,1]$.

\begin{center}
\includegraphics[scale=0.9]{courbeax2b.2}
\end{center}
\end{exemple}

\newpage
\section{Signes d'un polynôme du second degré}
\begin{prop}
\begin{itemize}
\item Si le polynôme $P$ n'a pas de racine alors 

\vspace{1ex}

\begin{tabular}[c]{|c|ccc|} \hline
$x$ & $-\infty$ &  & $+\infty$ \\ \hline
Signe de $P(x)$ &  & Signe de $a$ & \\ \hline
\end{tabular}

\item Si $x_1=x_2$ alors 
\vspace{1ex}

\begin{tabular}[c]{|c|ccccc|} \hline
$x$ & $-\infty$ &  & $-\frac{b}{2a}$ &  & $+\infty$ \\ \hline
Signe de $P(x)$ &  & Signe de $a$ & $0$ & Signe de $a$ & \\ \hline
\end{tabular}


\item Si $x_1 \neq x_2 \neq 0$ alors

\vspace{1ex}

\begin{tabular}[c]{|c|ccccccc|} \hline
$x$ & $-\infty$ &  & $x_{1}$ &  & $x_{2}$ &  & $+\infty$ \\ \hline
Signe de $P(x)$ &  & Signe de $a$ & $0$ & Signe de $-a$ & $0$ & Signe de $a$ & \\ \hline
\end{tabular}

\vspace{1ex}

\hspace*{1cm}où $x_{1}$ et $x_{2}$ sont les racines de $P$ et $x_{1}<x_{2}$
\end{itemize}
\end{prop}

%\begin{exemple}
%Pour étudier le signe de $2(x - 5)(x - \frac{5}{2})$ sur $\mathbb{R}$, nous pouvons utiliser un tableau de signe.

%\begin{tabular}[c]{|c|ccccccc|} \hline
%$x$ & $-\infty$ &  & $\frac{5}{2}$ &  & $5$ &  & $+\infty$ \\ \hline
%Signe de $2(x-5)(x-\frac{5}{2})$ &  & + & $0$ & - & $0$ & + & \\ \hline
%\end{tabular}

%\end{exemple}


\section{Représentation graphique, axe de symétrie}
\begin{prop}
Soit $\oij$ un repère du plan.

La représentation graphique de $P$ est l'image d'une translation de la parabole représentant la fonction $x\mapsto ax^{2}$.
L'axe de sym\'etrie d'une parabole repr\'esentative d'une fonction polyn\^ome du second degr\'e  du type $P(x)=ax^2+bx+c$ est $x=\dfrac{-b}{2a}$.
C'est donc une parabole de sommet $S\left(-\frac{b}{2a};P(-\frac{b}{2a})\right)$.
\end{prop}

\begin{center}
 \begin{tabular}{|c|c|c|c|}
   \hline
    & Deux racines distinctes  & Une racine double & Pas de racine  \\
   \hline
   Factorisation de $P(x)$ & $a(x-x_{1})(x-x_{2})$ & $a(x-x_{0})^{2}$ &
   pas de factorisation  \\
   \hline
   Équation $P(x)=0$ & $2$ solutions $x_{1}$ et $x_{2}$ & une solution $x_{0}$ & pas de solution\\ \hline
   \rule[0pt]{0pt}{2.2cm} \raisebox{1.1cm}[0pt][0pt]{$a>0$} & \includegraphics[scale=0.7]{CoursSecondDegFigures.1} & \includegraphics[scale=0.7]{CoursSecondDegFigures.2}
   & \includegraphics[scale=0.7]{CoursSecondDegFigures.3} \\
   \hline
   \rule[0pt]{0pt}{2.2cm} \raisebox{1.1cm}[0pt][0pt]{$a<0$} & \includegraphics[scale=0.7]{CoursSecondDegFigures.4} & \includegraphics[scale=0.7]{CoursSecondDegFigures.5} &  \includegraphics[scale=0.7]{CoursSecondDegFigures.6} \\
   \hline
  \end{tabular}
\end{center}

  \newpage



\section{Variations et extrema}

\begin{prop}
\begin{multicols}{2}
les variation d'un polynôme du second degré sont si $a<0$ :
\begin{itemize}
\item strictement croissante sur $]-\infty; \frac{-b}{2a}]$
\item strictement décroissante sur $[\frac{-b}{2a}; + \infty[$
\end{itemize}

les variation d'un polynôme du second degré sont si $a>0$ :
\begin{itemize}
\item strictement décroissante sur $]-\infty; \frac{-b}{2a}]$
\item strictement croissante sur $[\frac{-b}{2a}; + \infty[$
\end{itemize}
\end{multicols}
\end{prop}

\begin{prop}
La valeur de l'extrema d'un polynôme du second degré est atteinte en $x = \frac{-b}{2a}$ c'est un maximum si $a<0$ ou un minimum si $a>0$. 
\end{prop}
 
\begin{exemple}
 Considérons le polynôme du second degré $p(x) = -x^2 + 4x + 3$. Nous allons déterminer le maximum de ce polynôme.

Tout d'abord, nous pouvons remarquer que le coefficient du terme en $x^2$ est négatif, ce qui signifie que le graphe de la fonction $p$ est une parabole ouverte vers le bas, et donc que $p$ admet un maximum.\\

La valeur du maximum de $p$ est atteinte en $x = \frac{-b}{2a} = \frac{-4}{-2} = 2$. Nous avons donc $p(2) = -(2)^2 + 4(2) + 3 = 7$.

Ainsi, le maximum de $p$ est $7$, atteint en $x=2$. La courbe de la fonction $p$ est une parabole ouverte vers le bas, qui passe par le point $(2,7)$ est le sommet de cette parabole, qui est le point le plus élevé de la courbe.

%\begin{center}
% 
%\begin{tikzpicture}
%   \tkzTabInit{$x$ / 1 , Variation de $p$ / 1.5}{$-\infty$, $2$, $+\infty$}
%   \tkzTabVar{-/ $-\infty$ , +/ 7, -/ $-\infty$}
%\end{tikzpicture}
%
%
%\end{center}

%


Représentions graphique de $p$ sur $[-0,7;4,7]$.

\begin{center}
\includegraphics[scale=0.62]{courbeax2b.8}
\end{center}
\end{exemple}


\section{Forme factorisée et forme développée d'un polynôme du troisième degré}

\begin{defn}
On appelle \gap*{fonction polynôme du troisième degré} toute fonction $f$ définie sur $\mathbb{R}$ et qui s'écrit $f(x)=ax^3+bx^2+cx+d$ où $a$, $b$, $c$ et $d$ sont des réels fixés et $a\neq 0$.\\

Un polynôme du troisième degré peut être écrit sous \gap*{forme factorisée} sous la forme $a(x - x_1)(x - x_2)(x-x_3)$, où $x_1$, $x_2$ et $x_3$ sont les \gap*{racines du polynôme}
\end{defn}

\section{Les polynômes du troisième degré sous forme factorisée }
\begin{exemple}

Représentation graphique de $m$ définie par $m(x)=\dfrac{-2}{10}x(x-3)(x+2)$ sur $[-3;3,5]$.

\begin{center}
\includegraphics[scale=0.7]{courbeax2b.7}
\end{center}

\end{exemple}

\section{Les polynômes du troisième degré du type $x\mapsto ax^3$ et $x\mapsto ax^3+b$}

\begin{exemple}

Représentation graphique de $h$ définie par $h(x)=2x^3$ sur $[-3;3]$.

\begin{center}
\includegraphics[scale=0.7]{courbeax2b.4}
\end{center}

\end{exemple}

\begin{exemple}

Représentation graphique de $q$ définie par $q(x)=\dfrac{-1}{10}x^3$ sur $[-3;3]$.

\begin{center}
\includegraphics[scale=0.9]{courbeax2b.5}
\end{center}

\end{exemple}

\begin{exemple}

Représentation graphique de $n$ définie par $n(x)=\dfrac{-1}{10}x^3+1$ sur $[-3;3]$.

\begin{center}
\includegraphics[scale=0.9]{courbeax2b.6}
\end{center}

\end{exemple}


\section{Equation du type $x^3=c$, avec $c$ strictement positif}

\begin{thm}
Soit $c$ un nombre strictement positif alors l'équation du type $x^3=c$ admet une \gap*{unique solution} appelé \textit{racine cubique de $c$} et souvent noté $\sqrt[3]{c}$ (ou encore $c^{\frac{1}{3}}$).  
\end{thm}

\begin{exemple}
Pour résoudre une équation de la forme $x^3=c$ avec $c$ strictement positif à l'aide d'une calculatrice, nous pouvons utiliser la fonction racine cubique. Cette fonction est généralement notée $\sqrt[3]{\cdot}$ ou $\text{cubeRoot}(\cdot)$ sur les calculatrices.

Prenons l'exemple suivant : $x^3=125$. Nous allons utiliser la fonction racine cubique pour trouver la valeur de $x$.

Sur une calculatrice, nous pouvons entrer l'équation $x^3=125$ de la manière suivante :

\begin{enumerate}
\item Appuyez sur la touche de la racine cubique $\sqrt[3]{\cdot}$.
\item Appuyez sur la touche $125$.
\item Appuyez sur la touche $EXE$.
\end{enumerate}

Cela donnera la solution $x=5$.

Ainsi, la solution de l'équation $x^3=125$ est $x=5$.
\end{exemple}

\newpage

\section{Exercices}

\begin{exocor}
Développer, réduire et ordonner :
\begin{enumerate}
\begin{multicols}{2}
\item  $f(x)=(2x-3)(x+1)$
\item  $f(x)=(x-3)(4 x-1)$
\item  $f(x)=3(x-1)(x+2)$
\item $f(x)=-5(x-4)^{2}$
\end{multicols}
\end{enumerate}
\end{exocor}

\begin{exocor}
Représenter dans un repère orthogonal la fonction $f$ définie sur $\mathbb{R}$ par $f(x)=-2 x^{2}$.
\end{exocor}


\begin{exocor}
Résoudre dans $\mathbb{R}$ les équations suivantes:

\begin{enumerate}
\begin{multicols}{3}
\item  $x^{2}=5$
\item  $x^{2}=36$
\item  $x^{2}=7$
\item  $x^{2}=18$
\item  $x^{2}=0$
\item  $x^{2}=-27$
\end{multicols}
\end{enumerate}
\end{exocor}



\begin{exocor}
Soit $f$ la fonction définie sur $\mathbb{R}$ par $f(x)=\frac{1}{3} x^{2}$.

\begin{enumerate}
\item Dresser le tableau de variation de $f$.
\item Résoudre algébriquement l'équation $f(x)=3$.
\end{enumerate}
\end{exocor}

\begin{exocor}
Résoudre dans $\mathbb{R}$ les équations suivantes. On donnera la valeur exacte puis une valeur approchée à 0,01 près des solutions.

\begin{enumerate}
\begin{multicols}{3}
\item $x^{2}=100$
\item $x^{2}=64$
\item $2 x^{2}+1=5$
\item $x^{2}-4=0$
\item $2 x^{2}=x^{2}+1$
\item $x^{2}-4 x^{2}=0$
\end{multicols}
\end{enumerate}
\end{exocor}

\begin{exocor}
Représenter dans un repère orthogonal la fonction $f$ définie sur $\mathbb{R}$ par $f(x)=\frac{1}{3} x^{2}-2$.
\end{exocor}

\begin{exocor}
Résoudre dans $\mathbb{R}$ les équations suivantes :
\begin{enumerate}

\begin{multicols}{2}
\item $4 x^{2}-16=0$;
\item $-2 x^{2}+8=0$;
\item $3 x^{2}+7=1$
\item $7 x^{2}-5=10$
\end{multicols}
\end{enumerate}
\end{exocor}

\begin{exocor}
On considère les fonctions $f$ et $g$ dont les représentations graphiques sont les paraboles données ci-dessous. Pour chacune d'elle, déterminer son expression puis dresser son tableau de variation.

\begin{center}
\includegraphics[scale=1]{courbeax2b.11}
\end{center}

\end{exocor}

\begin{exocor}
Soit $f$ la fonction définie sur $\mathbb{R}$ par $f(x)=x^{2}-3$ et $\mathcal{C}$ sa courbe représentative. Indiquer si les affirmations suivantes sont vraies ou fausses. Justifier

\begin{enumerate}
\item Le minimum de $f$ est 0 .
\item Le sommet de $\mathcal{C}$ a pour coordonnées $S(0 ;-3)$.
\item L'axe de symétrie de $\mathcal{C}$ a pour équation $x=-3$.
\end{enumerate}
\end{exocor}


\begin{exocor}
Pour chaque polynôme du second degré, déterminer les racines :
\begin{enumerate}
\begin{multicols}{2}
\item $f(x)=2(x-3)(x+1)$
\item $f(x)=-5(x-3)(x-1)$
\item $f(x)=3(x-6)(x+2,5)$
\item $f(x)=5(x+8)^{2}$
\end{multicols}
\end{enumerate}
\end{exocor}

\begin{exocor}
On considère le polynôme du second degré $f(x)=-3 x^{2}-3 x+6$.

\begin{enumerate}
\item Vérifier que 1 et -2 sont des racines de $f$. 
\item En déduire une factorisation de $f(x)$.
\item Même question avec $f(x)=4 x^{2}-4$ et les racines -1 et 1 .
\item Même question avec $f(x)=-2 x^{2}+8 x-8$ et 2 pour unique racine (racine double).
\end{enumerate}
\end{exocor}

\begin{exocor}
On considère le polynôme du second degré $f(x)=3 x^{2}-15 x+18$.
\begin{enumerate}
\item Vérifier que $\alpha=2$ est une racine de $f$.
\item Déterminer la seconde racine de $f$.
\end{enumerate}
\end{exocor}

\begin{exocor}
On considère le polynôme du second degré $f(x)=3 x^{2}+9x-30$.
\begin{enumerate}
\item Vérifier que $\alpha=-5$ est une racine de $f$.
\item Déterminer la seconde racine de $f$.
\end{enumerate}
\end{exocor}

\begin{exocor}
Soit $f$ la fonction polynôme de second degré définie sur $\mathbb{R}$ par $f(x)=2 x^{2}-6 x-20$.

\begin{enumerate}
\item Vérifier que $2 x^{2}-6 x-20=2(x+2)(x-5)$.
\item En déduire les racines de $f$ et l'axe de symétrie de la courbe représentative de $f$.
\item Déterminer les coordonnées du sommet de la parabole, courbe représentative de $f$.
\end{enumerate}
\end{exocor}

\begin{exocor}
Soit $f$ la fonction polynôme du second degré définie sur $\mathbb{R}$ par $f(x)=-3 x^{2}-3 x+6$.

\begin{enumerate}
\item Vérifier que $f(x)=-3(x-1)(x+2)$ pour tout réel $x$.
\item En déduire les racines de $f$.
\item Dresser le tableau de variation de $f$ sur $\mathbb{R}$.
\end{enumerate}

\end{exocor}

\begin{exocor}
On considère la fonction $f$ du second degré définie par $f(x)=x^{2}-58 x-183$ sur $\mathbb{R}$.

\begin{enumerate}
\item  Vérifier que $f(x)=(x-61)(x+3)$ pour tout réel $x$.
\item  En déduire les racines de $f$.
\item  Dresser le tableau de variation de $f$ sur $\mathbb{R}$.
\item  En déduire le minimum de $f$ sur $\mathbb{R}$ et la valeur en laquelle il est atteint. 

\end{enumerate}
\end{exocor}

\begin{exocor}
On considère la fonction $f$ du second degré définie par $f(x)=3x^{2}+3 x-6$ sur $\mathbb{R}$.

\begin{enumerate}
\item  Vérifier que $f(x)=3(x-1)(x+2)$ pour tout réel $x$.
\item  En déduire les racines de $f$.
\item  Dresser le tableau de variation de $f$ sur $\mathbb{R}$.
\item  En déduire le minimum de $f$ sur $\mathbb{R}$ et la valeur en laquelle il est atteint. 

\end{enumerate}
\end{exocor}

\begin{exocor}
On considère la fonction $f$ définie sur l'intervalle $[-4 ; 4]$ par $f(x)=x^{2}-2 x-3$.
\begin{enumerate}
\item  Calculer l'image de -1 par $f$.
\item  Montrer que 3 est solution de l'équation $f(x)=0$.
\item  En utilisant les questions 1 et 2, donner une forme factorisée de $f(x)$.
\item  Dresser le tableau de signes de $f$ sur l'intervalle $[-4 ; 4]$.
\item Parmi les trois courbes suivantes, déterminer, en justifiant, celle qui représente graphiquement la fonction $f$.
\end{enumerate}

\begin{center}
\includegraphics[scale=0.8]{courbeax2b.12}
\end{center}

\end{exocor}

\begin{exocor}
Une usine agricole produit chaque mois entre 0 et 50 machines agricoles. On a modélisé le bénéfice de cette entreprise, exprimé en milliers d'euros, par la fonction $f$ définie pour tout $x \in[0 ; 50]$ par :

$$
f(x)=x^{3}-95 x^{2}+2450 x-10000
$$

\begin{enumerate}
\item Montrer que, pour tout $x \in[0 ; 50], f(x)=(x-5)(x-40)(x-50)$.
\item Étudier le signe de $f(x)$
\item En déduire le nombre de machines agricoles que l'entreprise doit produire pour réaliser des profits. 
\end{enumerate}

\end{exocor}


\begin{exocor}

\begin{enumerate}
\item Justifier que les fonctions suivantes sont des fonctions polynômes du troisième degré.

\begin{enumerate}
\item $f(x)=4 x^{3}-2 x+3$
\item $g(x)=-3(x+7)(x-2)(x-1)$
\item $h(x)=2(x-1)(x-1)^{2}$
\end{enumerate}

\item On considère la fonction polynôme de degré 3 définie pour tout $x \in \mathbb{R}$ par : $f(x)=x^{3}-3 x^{2}+4$.

\begin{enumerate}
\item Montrer que 2 et $-1$ sont deux racines de $f$.
\item Que peut-on en déduire graphiquement ?
\end{enumerate}

\end{enumerate}

\end{exocor}

\begin{exocor}
Résoudre dans $\mathbb{R}$ les équations suivantes. On donnera la valeur exacte puis une valeur approchée à 0,01 près des solutions.

\begin{enumerate}
\begin{multicols}{4}
\item $x^{3}=1000$
\item $x^{3}=64$
\item $2 x^{3}+11=5$
\item $x^{3}-4=0$
\item $2 x^{3}=x^{3}+1$
\item $x^{3}-4 x^{2}=0$
\item $x^{3}+3 x^{2}=3 x^{2}-1$ 
\item $3 x^{3}+3=2 x^{3}+4$
\end{multicols}
\end{enumerate}
\end{exocor}

\begin{exocor}
Soit $f$ la fonction définie sur $\mathbb{R}$ par :

$$f(x)=-3(x-8)(x+5)(x-3)$$

\begin{enumerate}
	\item Dresser le tableau de signes de $f$ sur $\mathbb{R}$.
	\item En déduire les solutions de l'inéquation $f(x)<0$.	
	\item Résoudre dans $\mathbb{R}$ les inéquations suivantes :
		\begin{enumerate}
			\item $2(x-1)(x+4)(x-3) \geqslant 0$
			\item $-(x+1)^{2}(x+2) \leqslant 0$
			\item $-2 x(x+5)(x-2)>0$  
			\item $(x-1)^{3} \leqslant 0$
		\end{enumerate}
\end{enumerate}
\end{exocor}

\begin{exocor}
L'objectif de l'exercice est de trouver le maximum de la fonction $f$ définie sur l'intervalle $[200 ; 400]$ par

$$f(x)=-0,01 x^{3}+4 x^{2}$$

\begin{enumerate}
	\item On admet que la fonction $f$ est dérivable sur $[200 ; 400]$ et on note $f^{\prime}$ sa dérivée.
	\item Calculer $f^{\prime}(x)$ et montrer que $f^{\prime}(x)=x(-0,03 x+8)$.
	\item Donner le tableau de signe de la fonction dérivée $f^{\prime}$ sur l'intervalle $[200 ; 400]$.
	\item En déduire le tableau de variation de la fonction $f$ sur l'intervalle $[200 ; 400]$.
	\item Quel est le maximum de cette fonction sur l'intervalle $[200 ; 400]$ ? En quelle valeur est-il atteint?
\end{enumerate}

\end{exocor}

\begin{exocor}
Soit $f$ la fonction définie sur l'intervalle $[-5 ; 5]$ par :

$$f(x)=x^{3}-3 x^{2}-24 x+8$$

\begin{enumerate}
	\item Calculer $f^{\prime}(x)$, où $f^{\prime}$ désigne la fonction dérivée de $f$ sur l'intervalle $[-5 ; 5]$.
	\item Vérifier que pour tout $x \in[-5 ; 5], f^{\prime}(x)=3(x-4)(x+2)$.
	\item Étudier le signe de $f^{\prime}(x)$ sur l'intervalle $[-5 ; 5]$.
	\item En déduire les variations de $f$ sur l'intervalle $[-5 ; 5]$.
	\item Déterminer la valeur de $x$ pour laquelle la fonction $f$ admet un maximum sur l'intervalle $[-5 ; 5]$ et en préciser la valeur.
\end{enumerate}

\end{exocor}

\newpage

\begin{exocor}
On a représenté graphiquement ci-dessous une fonction $f$ polynôme du troisième degré sur un intervalle et telle que

$$f(x)=a(x-x_1)(x-x_2)(x-x_3)$$

\begin{center}
\includegraphics[scale=0.8]{courbeax2b.9}
\end{center}
\begin{enumerate}
\item Quelles sont les valeurs de $x_{1}, x_{2}$ et $x_{3}$ ?
\item Déterminer la valeur de $a$ sachant que $f(0)=1$.
\item En déduire l'expression de $f(x)$.
\end{enumerate}

\end{exocor}

\begin{exocor}
On a représenté graphiquement ci-dessous une fonction $g$ polynôme du troisième degré sur un intervalle et telle que

$$g(x)=a(x-x_1)(x-x_2)(x-x_3)$$

\begin{center}
\includegraphics[scale=0.8]{courbeax2b.10}
\end{center}


\begin{enumerate}
\item Quelles sont les valeurs de $x_{1}, x_{2}$ et $x_{3}$ ?
\item Déterminer la valeur de $a$ sachant que $g(1)=-2$.
\item En déduire l'expression de $g(x)$.
\item Déterminer l'expression de $g'(x)$.
\item Montrer que $g'(x)=3x(x^2-3)$
\item Dresser le tableau de signe de $g'(x)$ et le tableau de variation de $g$ sur $[-2;2]$.
\end{enumerate}

\end{exocor}


\end{document}
